\chapter{Conclusions and Future Work}

\section{Conclusions}

This dissertation set out to design, \emph{in silico}, high-quality \acrshort{ssdna} capture
probes for \acrshort{gfet} biosensors targeting a panel of clinically relevant bacterial
pathogens. The main outcome is a single, reproducible computational pipeline that turns a
target gene into a small set of ranked, biophysically validated probe candidates, with every
design decision grounded in the literature.

The pipeline integrates, in one automated workflow, the stages that are usually performed
separately or by hand: retrieval and curation of sequences from \acrshort{ncbi}, automatic
length clustering, multiple sequence alignment and \acrshort{ppi}-based conservation scoring,
nearest-neighbour thermodynamic screening, secondary-structure evaluation with a continuous
No-fold score, and a final three-dimensional validation with Boltz-2. Two design choices give
the method its generality and rigour. First, thresholds are resolved through a layered,
\emph{organism-aware} scheme (gene\,$\rightarrow$\,species\,$\rightarrow$\,type\,$\rightarrow$\,
global), so that pathogens with very different genome composition are judged by fair criteria;
this was shown to recover candidates --- such as AT-rich probes --- that a single global rule
would wrongly discard. Second, candidates are prioritised by a transparent quality score
combining conservation and the No-fold score, which concentrates the expensive 3D validation on
the most promising sequences.

Applied to the six-gene panel, the pipeline reduced $8455$ candidate windows to $2943$ that
passed screening and identified a shortlist of top candidates whose high sequence quality was
independently corroborated by the highest 3D confidence and \acrshort{plddt}. Alignment-free
diversity and rarefaction analyses characterised each target and justified the sampling depth
in a principled, non-arbitrary way. An existing experimental probe panel was folded into the
same workflow and scored with identical metrics, enabling a direct comparison, and a curated
$150$-probe dataset of the highest-quality candidates per gene was produced.

Overall, the work delivers a solid, extensible foundation for the computational stage of
\acrshort{gfet} biosensor development: it produces prioritised, structurally sound probe
candidates and generalises to new genes and organisms with minimal configuration.

\section{Prospects for future work}

Several directions follow naturally from this work:

\begin{itemize}
  \item \textbf{Experimental validation.} The ultimate test of the designed probes is their
        performance on fabricated \acrshort{gfet} devices; the ranked shortlist provides the
        immediate candidates for functionalisation and characterisation.
  \item \textbf{Specificity screening.} Adding an automated specificity check (for example a
        BLAST step against non-target genomes and the host) would complement the current
        conservation-driven, sensitivity-oriented selection.
  \item \textbf{Comparison of 3D methods.} The Boltz-2 predictions can be benchmarked against an
        independent molecular-docking pipeline, to assess the agreement of the two structural
        views of the probe--target interaction.
  \item \textbf{Data-driven modelling.} The curated feature dataset opens the way to a machine-
        learning model that predicts probe quality directly, potentially accelerating or
        refining the ranking.
  \item \textbf{Broader panels.} As the framework already supports viral, fungal, protozoan and
        host profiles, extending the panel to further genes and non-bacterial pathogens is
        largely a matter of configuration.
\end{itemize}
